\documentclass[10pt,a4paper]{amsart}
\usepackage[margin=19mm]{geometry}
\usepackage{amsmath,amssymb,mathtools}
\usepackage[scr=rsfs,cal=euler]{mathalfa}
\usepackage{xfrac}
\usepackage[unicode,hidelinks,bookmarks=false]{hyperref}
\newcommand{\ie}{i.e.\ }
\newcommand{\cf}{cf.\ }
\newcommand{\resp}{respectively\ }
\newcommand{\Subsection}{\subsection}
\providecommand{\nobreakdash}{\nobreak\hbox{-}}
\setlength{\emergencystretch}{2em}
\multlinegap0pt
\usepackage{amssymb}
\usepackage{mathtools}
\newtheorem{theorem}{Theorem}
\title{Maximal discrete subgroups in unitary groups \\ of operator algebras}
\author{Vadim Alekseev, Andreas Thom}
\date{Source: arXiv:2206.06704v1 (CC BY 4.0)}
\begin{document}
\maketitle

\noindent\emph{Source and licence.} Maximal discrete subgroups in unitary groups \\ of operator algebras. Version arXiv:2206.06704v1 (CC BY 4.0), distributed by arXiv under the Creative Commons Attribution 4.0 International licence, http://creativecommons.org/licenses/by/4.0/, as recorded in the arXiv licence metadata for this exact version and shown on the abstract page https://arxiv.org/abs/2206.06704v1. Full licence notice: \texttt{licenses/arxiv-2206.06704-CC-BY-4.0.txt}.\par\smallskip

\noindent\textit{Editorial scope.} Counted unit: the article's theorem cstar (source0.tex lines 236--240). The article's complete original proof of it is delivered with it (source0.tex lines 241--253, one \texttt{proof} environment of 13 lines). No mathematics is written, repaired or re-derived: the statement and the proof are the article's own lines, in the article's own notation, and no retained line is substituted or re-flowed.  No further source-local context is needed: every cross-reference of every retained line resolves inside the two delivered spans.  Nothing was omitted from any retained span, so the package carries no omission record.  No retained line contains a citation, so the wrapper carries no bibliography and no bibliography entry.  No retained line contains a figure environment, an image-inclusion command or drawing code, so no image is delivered and the package declares no asset dependency.  Editorial formatting only, in addition: an AMS \texttt{amsart} preamble that restates the article's own declarations and macros used by the retained lines, namely the environments \texttt{theorem} on separate counters, together with the standard packages those lines need.  The retained environments print in this document as Theorem 1 in order of appearance, whereas the article prints the same items with the numbering of its own full document.  The complete native source of the article as distributed in the arXiv e-print for this exact version is bundled unchanged as \texttt{sources/126912/source0.tex}.
\begin{theorem}\label{cstar}
Let $A$ be a unital $C^*$-algebra and $\Gamma$ be a discrete subgroup of its unitary group ${\rm U}(A)$. Then, the subgroup
$$\Gamma_{1/2} = \langle g \in \Gamma \mid \|1-g\|<1/2 \rangle$$
is abelian and normal in $\Gamma$.
\end{theorem}

\begin{proof}
We set $\ell(g)=\|1-g\|$ and note that
\begin{eqnarray*} \ell([g,h]) &=& \|1-ghg^*h^*\| = \|hg-gh\| \\
&=& \|(1-h)(1-g)-(1-g)(1-h)\| \leq 2 \|1-g\|\|1-h\| = 2 \ell(g) \ell(h).
\end{eqnarray*}
We define $\Gamma_{t} := \langle g \in \Gamma \mid \ell(g)<t \rangle$. 
%It follows from the standard commutator identities and the inequality above that $[\Gamma_t,\Gamma_{t'}] \subset \Gamma_{4tt'}.$ In particular, $\Gamma_{1/5}$ is nilpotent, since $n$-fold iterated commutators will be contained in $\Gamma_{(4/5)^n}$ which is eventually trivial. We will now show that $\Gamma_{1/5}$ is abelian. Consider 
Fix $\varepsilon>0$ and set
$$\delta = \inf\{\ell(g) \mid g \in \Gamma_{1/2-\varepsilon} \mbox{ is not central in } \Gamma_{1/2-\varepsilon}\}.$$
Suppose that this infimum is positive and finite, i.e.\ that the set of non-central elements is bounded away from $1$ and non-empty. If $g$ is non-central in $\Gamma_{1/2-\varepsilon}$, then it cannot commute with the entire generating set of $\Gamma_{1/2-\varepsilon}$. Thus, we can consider $g,h \in \Gamma_{1/2-\varepsilon}$ with $\ell(g) < (1+2\varepsilon)\delta$, $\ell(h)<1/2 - \varepsilon$ and $[g,h] \neq 1$. We compute $$\ell([g,h]) \leq 2 (1+2\varepsilon)\delta (1/2-\varepsilon)<\delta$$ and hence $[g,h]$ is central in $\Gamma_{1/2-\varepsilon}$, in particular, it commutes with $g$ and $h$. Thus, we obtain a unitary representation of the Heisenberg group
$$H(\mathbb Z) = \langle x,y,z \mid [x,y]=z, [x,z]=[y,z]=1 \rangle.$$ But for any unitary representation of the Heisenberg group for which $z$ acts non-trivially, both generators are far away from the identity. Indeed, irreducible unitary representations are parametrised by $\theta \in S^1$ and either lead to the unique irreducible representation of the non-commutative torus $A_{\theta}$ or they are finite-dimensional and completely understood. In any case, the generators satisfy $\|1-g\| \geq \sqrt{3}$ and $\|1-h\| \geq \sqrt{3}$ unless they commute. This implies that $\Gamma_{1/2-\varepsilon}$ is abelian. Since $\varepsilon>0$ was arbitrary and $$\Gamma_{1/2} = \bigcup_{\varepsilon>0} \Gamma_{1/2-\varepsilon},$$
this finishes the proof.
\end{proof}

\end{document}
